\begin{abstract}
Held-out-form benchmarks for neuro-symbolic expression encoders are meant to test whether a model recognises algebraic \emph{structure} or merely which \emph{tokens} appear. We give a four-tier account of what such a benchmark can be solved by (variable identity, operator and arity inventory, order and local tree shape, and genuine composition), release a training-free auditor that measures the lower tiers, and state a formal criterion under which held-out-form identification constitutes evidence of \emph{compositional} generalization. The third tier is the one usually skipped, and it carries most of the apparent headroom: on EQNET \texttt{poly8} a zero-parameter bag over parent--child and subtree-size statistics reaches $0.411$ where the operator bag reaches $0.186$, so clearing an operator bag establishes far less than it appears to.
On \texttt{poly8} (1102 equivalence classes with a shared variable pool and \emph{zero} classes having a unique operator bag), trained Tree-LSTMs reach unseen-class $k$-NN precision of $0.896$ at $K{=}500$ against the strongest non-learned baseline's $0.518$, a margin that \emph{grows} with $K$. To hold the lower tiers fixed we pair each held-out form with a \emph{twin} of identical token and operator multisets differing only in how they combine. A rotation twin pins inventory at chance but still leaks shape (our tree-local bag scores $0.861$ on it), so we build a \emph{shape-matched} twin by swapping siblings under a non-commutative operator, which leaves that bag exactly invariant. With inventory and shape both pinned, the untrained encoder falls from $0.925$ to $0.788$ while the trained one holds at $0.994$: the trained--untrained gap \emph{widens} to $0.206$. That an untrained recursive encoder scores $0.788$ at all is why we report such baselines as \emph{non-learned} rather than ``structure-free''.
The framework also locates the flaw the field worries about. A per-symbol tokenization convention leaves $84\%$ of AI~Feynman equations with a variable set no other equation shares, so on the resulting suite a zero-parameter bag over variable tokens alone is a complete lookup solution by construction: it scores $1.000$ against a trained Tree-LSTM's $0.972$; but the canonical published equivalence-class benchmarks reuse a shared variable pool and are immune at tier~1. Tier~2 is where the exposure is, almost everywhere.
We hold our own results to the same standard, and it costs us: both purpose-built ``leakage-proof'' suites turn out to be solvable at tier~2, and the protocol-sensitivity contrast shrinks from $0.44$ to $0.13$ under corrected provenance. The shared-variable Feynman redesign survives ($0.932$ versus $0.564$ variable-blind, $0.604$ untrained), and a leave-one-schema-out experiment holding the held-out tree byte-identical shows library coverage is a causal part of what remains.
\end{abstract}
